\subsection{Opponent Shaping Leads to the Spontaneous Rediscovery of Game-Theoretic Strategies~\citep{lu2022modelfreeopponentshaping, foerster2018learningopponentlearningawareness}}
\label{sec:opponentshaping}

The final anecdote in this section steps back from the large, highly engineered systems discussed previously to comparatively small-scale multi-agent learning experiments at the intersection of game theory and reinforcement learning, asking whether similarly surprising strategic phenomena appear even in a much simpler repeated-game setting.
Researchers used RL algorithms to study repeated games~\citep{aumannrepeated59}---situations where the same players face each other over and over again.
In a single game, it often makes sense to be selfish; however, in a repeated game, players must weigh their immediate gains against the risk of future retaliation by the opponent.
% The two learned solutions in this anecdote displayed complex social behaviors.

Learning in multi-agent environments~\citep{lanctot2017unifiedgametheoreticapproachmultiagent} is complex because the presence of other learning agents makes the environment non-stationary~\citep{amato2025initialintroductioncooperativemultiagent}: as one agent's strategy evolves, it changes the optimal strategy for other agents in the environment. Typically, agents are trained to improve their own policies myopically, without considering how their actions will influence the future learning behavior of others.
In an attempt to address this shortcoming, researchers explored algorithms designed to enable agents in multi-agent interactions to anticipate and influence their opponents' learning processes. Specifically, in the paper introducing LOLA (Learning with Opponent-Learning Awareness;~\citealp{foerster2018learningopponentlearningawareness}), agents explicitly accounted for how their actions affected their opponents' policy updates by having an explicit model of their opponent.
While researchers initially observed that independently trained agents in the Iterated Prisoner's Dilemma continuously defected, they were surprised when the LOLA agents instead spontaneously (re)discovered the cooperative strategy known as tit-for-tat~\citep{Axelrod1981gametheory}, where the agents first cooperate and only defect if the other player betrays them. %, even when the individual agents were incentivized to be selfish in the short term.

Similarly, in their follow-up work on M-FOS (Model-Free Opponent Shaping;~\citealp{lu2022modelfreeopponentshaping}), the researchers introduced a framework enabling agents to shape opponents' learning behavior without an explicit model of the opponent's behavior (hence, model-free). Remarkably, the M-FOS-trained agents independently discovered sophisticated strategies such as zero-determinant extortion~\citep{Stewart2013zerodeterminant, Hao2015ZeroDet, Press2012Zerodeterminant}, which had only recently been identified by human theorists within the last fifteen years.
The emergence of such complex strategies from model-free RL methods was unexpected and notable to the researchers.
